\subsection{Umbral Recurrences and Variants}

\begin{theorem} \label{thm:associated-sequence-recurrence-and-transfer} \hlt{Recurrence and Transfer Formulas}\\
Let $(p_n)$ be associated with the delta series $f$, and set $u(t)\defeq f(t)/t\in\mathscr{F}^{\times}$.\\
Then $p_{n+1}=x[f'(t)]^{-1}p_n$, and, for $n\geq1$, $p_n=f'(t)u(t)^{-n-1}x^n=xu(t)^{-n}x^{n-1}$.
\begin{proof} By the chain rule for umbral shifts with coordinate $t$, $\theta_t=(\partial_t f)\theta_f=f'(t)\theta_f$.\\
Since $\theta_t=x$ and $\theta_f p_n=p_{n+1}$, $p_{n+1}=x[f'(t)]^{-1}p_n$. The transfer identity follows by applying this recurrence inductively after writing $f=tu$, hence $p_n=f'(t)u(t)^{-n-1}x^n=xu(t)^{-n}x^{n-1}$.\\
\hlt{Structural Commentary}: The recurrence generates $(p_n)$ from $f'(t)^{-1}$.\\
\hlt{Key Insight}: Transfer formulas convert the difficult inverse coordinate problem into formal operator algebra.
\end{proof}
\end{theorem}

\begin{corollary} \label{cor:sheffer-shift-formula} \hlt{Sheffer Shift Formula}\\
If $(s_n)$ is Sheffer for $(g,f)$, then $\theta_{g,f}=g(t)^{-1}\theta_fg(t) =\left[x-\frac{g'(t)}{g(t)}\right]\frac{1}{f'(t)}$, and hence $s_{n+1}=\left(x-\frac{g'}{g}\right)(f')^{-1}s_n$.
\begin{proof}
Since $g(t)s_n$ is associated with $f$, $\theta_{g,f}=g^{-1}\theta_fg$. By Theorem \ref{thm:associated-sequence-recurrence-and-transfer}, $\theta_f=x(f')^{-1}$.\\
Using the Pincherle commutator $g'(t)=g(t)x-xg(t)$, we get $g^{-1}x(f')^{-1}g =g^{-1}xg(f')^{-1} =\left(x-\frac{g'}{g}\right)(f')^{-1}$.\\
Applying this operator to $s_n$ gives the recurrence.\\
\hlt{Structural Commentary}: The Sheffer shift is a conjugated associated
shift.\\
\hlt{Key Insight}: The logarithmic derivative $g'/g$ is the correction term.
\end{proof}
\end{corollary}

\begin{definition} \label{def:stirling-and-exponential-polynomials} \hlt{Stirling and Exponential Polynomials}
\begin{enumerate}[label=(\roman*)]
\setlength{\itemsep}{0pt}
\item \hlt{Stirling numbers}: $S(n,k)$ are defined by $x^n=\sum\limits_{k=0}^{n}S(n,k)(x)_k$.
\item \hlt{Exponential polynomials}: define $\phi_n(x)\defeq\sum\limits_{k=0}^{n}S(n,k)x^k$.
\item \hlt{Boundary convention}: $S(n,-1)=0$ and $S(n,k)=0$ for $k>n$.
\end{enumerate}
\end{definition}

\begin{example} \label{ex:umbral-recurrences-stirling} \hlt{Lower Factorials and Stirling Numbers}\\
For $f(t)=e^t-1$, $f'(t)=e^t$, hence $p_{n+1}(x)=xe^{-t}p_n(x)=xp_n(x-1)$. With $p_0=1$, this gives $p_n=(x)_n$.\\
For $f(t)=\log(1+t)$, $\overline f(t)=e^t-1$, and the associated sequence is $(\phi_n)$.\\
The inverse umbral relation is $x^n=\sum\limits_{k=0}^{n}S(n,k)(x)_k$, and $\phi_n(x)=\sum\limits_{k=0}^{n}S(n,k)x^k$.\\
The recurrence formula gives $\phi_{n+1}(x)=x(1+t)\phi_n(x)=x\bigl(\phi_n(x)+\phi_n'(x)\bigr)$,\\
so coefficient comparison yields $S(n+1,k)=S(n,k-1)+kS(n,k)$.
\end{example}

\begin{definition} \label{def:q-and-logarithmic-umbral-variants} \hlt{\texorpdfstring{$q$}{q}- and Logarithmic Umbral Variants}
\begin{enumerate}[label=(\roman*)]
\setlength{\itemsep}{0pt}
\item \hlt{Gaussian coefficient}: for $0\leq k\leq n$,
$\binom{n}{k}_q\defeq
\prod\limits_{j=1}^{n}(1-q^j)/
\left(\prod\limits_{j=1}^{k}(1-q^j)
\prod\limits_{j=1}^{n-k}(1-q^j)\right)$.
\item \hlt{\texorpdfstring{$q$}{q}-umbral calculus}: the deformation $\binom{n}{k}\leadsto\binom{n}{k}_q$.
\item \hlt{Logarithmic umbral calculus}: the logarithmic-type variant based on harmonic logarithms.
\end{enumerate}
\end{definition}

\begin{remark} \label{rem:umbral-calculus-further-directions} \hlt{Further Directions}\\
Classical umbral methods extend to many polynomial families and operations:
\begin{enumerate}[label=(\roman*)]
\setlength{\itemsep}{0pt}
\item \hlt{Families}: lower factorial, Stirling, exponential, Gould, central
factorial, Abel, Mittag--Leffler, Bessel, Bell, Hermite, Bernoulli, Euler,
Laguerre, Poisson--Charlier, actuarial, Meixner, Pidduck, Narumi, Boole,
Peters, Mahler, and Mott polynomials.
\item \hlt{Operations}: connection constants, duplication formulas, Lagrange
inversion, cross sequences, Steffensen sequences, operational formulas,
inverse relations, recurrence solving, and binomial convolution.
\item \hlt{Variants}: $q$-umbral and logarithmic umbral calculi replace or
modify the binomial structure through the objects in Definition
\ref{def:q-and-logarithmic-umbral-variants}.
\end{enumerate}
\end{remark}
